\chapter*{Abstract}
\addcontentsline{toc}{chapter}{Abstract}

A manipulated pressure reading may still look plausible. Different attacks leave different patterns over time, and missing measurements remove information that could help identify them. Combining specialist detectors creates a further problem: even accurate specialists can produce too many false alarms when used together. This thesis develops a hybrid detector to balance attack-family coverage with background errors.

General, Drift and Noise form the three main branches. General combines five tree-based components whose residual inputs come from a learned reference that excludes the target sensor group. Drift and Noise use temporal evidence with a declared three-hour decision delay. Their alarms pass through routing, feedback and verification, while General keeps its immediate decision path. For L-Town, General also receives 42 features calculated from actually received pressure history.

Evaluation uses simulated scenarios from Modena and L-Town. Pressure and flow missing probabilities are fixed at 0.50, and only observed pressure endpoints are scored. Ten model seeds and six evaluation sources per selected system separate variation in fitting from variation between sources. Mean pooled F1 is 0.8669 in Modena and 0.8950 in L-Town, with family macro F1 of 0.9045 and 0.9083. The respective clean-period false-positive rates are 0.0584\% and 0.0550\%. All five Modena family means exceed 0.80. Four exceed 0.89 in L-Town, where replay remains weaker at 0.7312.

The matched comparisons show a substantial reduction in Modena's false positives and higher L-Town family means after adding received history. Supplementary single-classifier controls put these gains in perspective. In both networks, the hybrid detects replay, drift and noise better but produces more clean false alarms. The simpler three-hour model remains competitive on L-Town pooled F1.
